\chapter{Ecuaciones diferenciales lineales}\label{ch:b1:diffeq}

Vistas por primera vez en el volumen anterior, las ecuaciones
diferenciales se tratan aquí con demostraciones completas y con mayor
generalidad: ecuaciones lineales de primer orden con coeficientes
variables (resueltas por completo con el método de \omterm{thm:b1:diffeq:voc}{variación de las
constantes}) y ecuaciones lineales de segundo orden con coeficientes
constantes, el modelo de las oscilaciones. Los dos casos exhiben la
misma estructura: \emph{solución general $=$ una solución particular
$+$ solución general de la ecuación homogénea}.

\section{Ecuaciones lineales de primer orden}

\begin{definition}\label{def:b1:diffeq:linear1}
Sean $I$ un intervalo y $a, b \colon I \to \R$ (o $\C$) continuas. La
ecuación
\[
(E)\colon\quad y' + a(x)\,y = b(x),
\]
en la incógnita, una función derivable $y \colon I \to \R$ (o $\C$),
es una \emph{ecuación diferencial lineal de primer
orden}\index{ecuación diferencial!lineal de primer orden}. La ecuación
$(H)\colon y' + a(x) y = 0$ es su ecuación \emph{homogénea}.
\end{definition}

\begin{theorem}[Resolución de la ecuación homogénea]\label{thm:b1:diffeq:homogeneous1}
Sea $A$ una primitiva de $a$ en $I$ (existe:
\cref{ch:b1:integration}). Las soluciones de $(H)$ en $I$ son
exactamente las funciones
\[
y(x) = \lambda\, \eu^{-A(x)}, \qquad \lambda \in \R \text{ (or } \C).
\]
\end{theorem}

\begin{proof}
Estas funciones son soluciones: $y' = -\lambda A' \eu^{-A} = -a y$.
Recíprocamente, sea $y$ solución de $(H)$ y póngase $z(x) = y(x)\, \eu^{A(x)}$. Entonces
\[
z' = y' \eu^{A} + y\, a\, \eu^{A} = (y' + ay)\, \eu^{A} = 0 ,
\]
luego $z$ es constante en el intervalo $I$, digamos $z = \lambda$:
$y = \lambda \eu^{-A}$. (Obsérvese la lógica: no se pierde ninguna
solución, porque \emph{toda} solución se ha escrito en la forma anunciada.)
\end{proof}

\begin{example}[Una ecuación homogénea de coeficiente variable]\label{ex:b1:diffeq:varcoeff}
Resuélvase $y' + (\cos x)\,y = 0$ en $\R$. Una primitiva de
$a(x) = \cos x$ es $A(x) = \sin x$, luego las soluciones son
\[
y(x) = \lambda\,\eu^{-\sin x}, \qquad \lambda \in \R .
\]
Dos lecturas. Toda solución es periódica (de período $2\pi$) y no se
anula nunca salvo si $\lambda = 0$ --- el signo de $\lambda$ es para
siempre el signo de $y$, pues una exponencial no puede cruzar el cero.
Y la solución con $y(0) = y_0$ es $y_0\eu^{-\sin x}$: exactamente una
curva de la familia por cada punto inicial, la imagen unidimensional
del~\cref{thm:b1:diffeq:voc}~(2).
\end{example}

\begin{theorem}[Variación de las constantes; problema de Cauchy]\label{thm:b1:diffeq:voc}
Con la notación anterior:
\begin{enumerate}
  \item Las soluciones de $(E)$ en $I$ son exactamente
        \[
        y(x) = \Bigl(\lambda + \int_{x_0}^{x} b(t)\, \eu^{A(t)} \dd t
        \Bigr)\, \eu^{-A(x)},
        \qquad \lambda \in \R,
        \]
        donde $x_0 \in I$ está fijado. Equivalentemente: solución
        general de $(H)$ más una solución particular de $(E)$.
  \item Para todos $x_0 \in I$ e $y_0$, el \emph{problema de
        Cauchy}\index{problema de Cauchy} «$(E)$ e $y(x_0) = y_0$»
        tiene exactamente una solución en $I$.
\end{enumerate}
\end{theorem}

\begin{proof}
(1) Siguiendo el método llamado \emph{variación de las
constantes}\index{variación de las constantes}, se buscan soluciones de
la forma $y = \mu(x)\, \eu^{-A(x)}$ con $\mu$ derivable --- no se
pierde generalidad, ya que toda función de $I$ se puede escribir así
($\mu = y\,\eu^{A}$). Sustituyendo,
\[
y' + ay = \mu' \eu^{-A} - \mu a \eu^{-A} + a\mu\eu^{-A}
= \mu'\, \eu^{-A},
\]
luego $y$ es solución de $(E)$ si y solo si $\mu'(x) = b(x)\,\eu^{A(x)}$,
si y solo si $\mu(x) = \lambda + \int_{x_0}^x b(t)\eu^{A(t)}\dd t$ para
alguna constante $\lambda$ (dos primitivas de una misma función
continua en un intervalo difieren en una constante).

(2) En la fórmula, $y(x_0) = \lambda\,\eu^{-A(x_0)}$: la condición
$y(x_0) = y_0$ determina $\lambda = y_0 \eu^{A(x_0)}$ de manera única.
\end{proof}

\begin{example}\label{ex:b1:diffeq:order1}
Resuélvase $y' + \dfrac{y}{x} = x^2$ en $I = \intoo{0}{+\infty}$. Aquí
$a(x) = \frac 1x$, $A(x) = \ln x$, $\eu^{-A(x)} = \frac 1x$. Soluciones
homogéneas: $\frac{\lambda}{x}$. \omterm{thm:b1:diffeq:voc}{Variación de las constantes}:
$\mu'(x) = x^2 \cdot x = x^3$, luego $\mu = \frac{x^4}{4} + \lambda$, y
\[
y(x) = \frac{x^3}{4} + \frac{\lambda}{x}, \qquad \lambda \in \R .
\]
Con la condición inicial $y(1) = 0$: $\lambda = -\frac14$.
\emph{Comprobación:} $y' + \frac yx = \frac{3x^2}{4} -
\frac{\lambda}{x^2} + \frac{x^2}{4} + \frac{\lambda}{x^2} = x^2$.
\end{example}

\begin{example}[Adivinar gana a integrar]\label{ex:b1:diffeq:guess}
Resuélvase $y' + 2x\,y = x$ en $\R$. La \omterm{thm:b1:diffeq:voc}{variación de las constantes}
funciona ($A = x^2$, $\mu' = x\,\eu^{x^2}$,
$\mu = \frac12\eu^{x^2} + \lambda$), pero observar que la
\emph{constante} $y_p = \frac12$ resuelve la ecuación
($0 + 2x\cdot\frac12 = x$) es más rápido. Con las soluciones
homogéneas $\lambda\,\eu^{-x^2}$:
\[
y(x) = \frac12 + \lambda\,\eu^{-x^2}, \qquad \lambda \in \R .
\]
Toda solución converge a $\frac12$ con enorme rapidez cuando
$x \to \pm\infty$: la solución particular constante es un
\emph{equilibrio} al que se unen todas las demás. La idea clave: antes
de lanzar el método general, dedíquense diez segundos a buscar una
solución particular evidente (constante, monomio, múltiplo del miembro
derecho); el teorema de estructura remata después el trabajo.
\end{example}

\begin{remark}[Los intervalos importan]
El teorema vive en un \emph{intervalo} donde $a$ y $b$ son continuas.
Para $y' + \frac yx = 0$ en $\R^*$, las soluciones son
$\frac{\lambda}{x}$ en $\intoo{0}{+\infty}$ y $\frac{\mu}{x}$ en
$\intoo{-\infty}{0}$ con constantes \emph{independientes}: no hay
ninguna razón para que una sola fórmula pegue las dos ramas a través de
la singularidad en $0$.
\end{remark}

\begin{example}[Un miembro derecho complejo, dos respuestas reales]\label{ex:b1:diffeq:complexrhs}
Resuélvanse $y' - y = \cos x$ e $y' - y = \sin x$ de una vez. Se
trabaja en $\C$ con el miembro derecho $\eu^{\iu x}$: probando
$y_p = c\,\eu^{\iu x}$ se obtiene $c(\iu - 1)\eu^{\iu x} = \eu^{\iu x}$, luego
\[
c = \frac1{\iu - 1} = \frac{-1 - \iu}2,
\qquad
y_p = -\frac{(1 + \iu)(\cos x + \iu\sin x)}2
= \frac{\sin x - \cos x}2 + \iu\,\frac{-\sin x - \cos x}2 .
\]
Como la ecuación tiene coeficientes reales, las partes real e
imaginaria se separan: $\frac{\sin x - \cos x}2$ resuelve
$y' - y = \cos x$, y $-\frac{\sin x + \cos x}2$ resuelve
$y' - y = \sin x$ (compruébese la primera: su derivada,
$\frac{\cos x + \sin x}2$, menos la función da $\cos x$). Una línea
compleja ha sustituido a dos pasadas de \omterm{thm:b1:diffeq:voc}{variación de las constantes}
--- la misma economía que el~\cref{met:b1:diffeq:particular}
sistematiza para el segundo orden, y un dividendo recurrente
del~\cref{ch:b1:complex}.
\end{example}

\section{Ecuaciones lineales de segundo orden con coeficientes constantes}

\begin{definition}\label{def:b1:diffeq:linear2}
Sean $a, b \in \R$ y $f \colon I \to \R$ continua. La ecuación
\[
(E)\colon\quad y'' + a\,y' + b\,y = f(x)
\]
es una \emph{ecuación lineal de segundo orden con coeficientes
constantes}\index{ecuación diferencial!lineal de segundo orden};
$(H)\colon y'' + ay' + by = 0$ es su ecuación homogénea, y
$\chi(r) = r^2 + ar + b$ su \emph{polinomio
característico}\index{polinomio característico}.
\end{definition}

\begin{theorem}[Soluciones homogéneas]\label{thm:b1:diffeq:homogeneous2}
Sea $\Delta = a^2 - 4b$ el discriminante de $\chi$. Las soluciones
reales de $(H)$ en $\R$ son:
\begin{enumerate}
  \item si $\Delta > 0$, con $r_1 \neq r_2$ las dos raíces reales:
        $\;y = \lambda\, \eu^{r_1 x} + \mu\, \eu^{r_2 x}$;
  \item si $\Delta = 0$, con $r_0$ la raíz doble:
        $\;y = (\lambda + \mu x)\, \eu^{r_0 x}$;
  \item si $\Delta < 0$, con raíces $\alpha \pm \iu\omega$
        ($\omega > 0$):
        $\;y = \eu^{\alpha x} (\lambda \cos\omega x + \mu
        \sin\omega x)$;
\end{enumerate}
en cada caso con $(\lambda, \mu)$ recorriendo $\R^2$.
\end{theorem}

\begin{proof}
Obsérvese primero que, para $r \in \C$, la función
$x \mapsto \eu^{rx}$ resuelve $(H)$ si y solo si $\chi(r) = 0$
(sustitúyase: $(r^2 + ar + b)\eu^{rx} = 0$). Por eso las exponenciales
son la primera conjetura natural: derivar actúa sobre $\eu^{rx}$ como
multiplicar por el número $r$, de modo que la ecuación diferencial se
convierte en la ecuación numérica $\chi(r) = 0$ --- todo el problema
analítico queda comprimido en hallar las raíces de una cuadrática.

El paso clave es un cambio de incógnita que rebaja el orden. Sea $r$
una raíz (posiblemente compleja) de $\chi$ y escríbase
$y = z\, \eu^{rx}$, lo que no pierde generalidad. Entonces
\[
y'' + ay' + by
= \bigl(z'' + (2r + a) z' + \chi(r) z\bigr)\eu^{rx}
= \bigl(z'' + (2r + a) z'\bigr)\eu^{rx},
\]
de modo que $(H)$ se convierte en la ecuación de \emph{primer orden}
$u' + (2r + a) u = 0$ para $u = z'$.

\emph{Caso $\Delta \neq 0$:} tómese $r = r_1$; entonces
$2r_1 + a = r_1 - r_2$ (pues $r_1 + r_2 = -a$). Por
el~\cref{thm:b1:diffeq:homogeneous1}, $z' = c\,\eu^{(r_2 - r_1)x}$ para
alguna constante $c$; integrando en $\R$,
$z = \mu\, \eu^{(r_2 - r_1)x} + \lambda$ con
$\mu = \frac{c}{r_2 - r_1}$, y por tanto
$y = z\,\eu^{r_1 x} = \lambda\,\eu^{r_1x} + \mu\,\eu^{r_2x}$.
Cuando $\Delta < 0$, las raíces son $\alpha \pm \iu\omega$ y las
soluciones complejas son $y = c_1\eu^{(\alpha+\iu\omega)x} +
c_2\eu^{(\alpha-\iu\omega)x}$ con $c_1, c_2 \in \C$. ¿Cuáles de ellas
toman valores reales? Como $\conj{\eu^{(\alpha+\iu\omega)x}} =
\eu^{(\alpha-\iu\omega)x}$, el \omterm{def:b1:complex:field}{conjugado} de $y$ es $\conj{c_2}\,
\eu^{(\alpha+\iu\omega)x} + \conj{c_1}\,\eu^{(\alpha-\iu\omega)x}$, y
que $y = \conj y$ para todo $x$ obliga a $c_2 = \conj{c_1}$ (las dos
exponenciales son linealmente independientes: evalúese en dos puntos, o
compárese en $x=0$ tras dividir por $\eu^{\alpha x}$). Escribiendo
$c_1 = \frac{\lambda - \iu\mu}2$ con $\lambda, \mu$ reales:
\[
y = 2\,\Re\Bigl(c_1\,\eu^{\alpha x}(\cos\omega x +
\iu\sin\omega x)\Bigr)
= \eu^{\alpha x}\bigl(\lambda\cos\omega x + \mu\sin\omega
x\bigr),
\]
y, recíprocamente, toda función así es solución (parte real de una
solución compleja de una ecuación real): el espacio de soluciones
reales es el anunciado.

\emph{Caso $\Delta = 0$:} $r = r_0$, $2r_0 + a = 0$, luego $z'' = 0$:
$z = \lambda + \mu x$ e $y = (\lambda + \mu x)\eu^{r_0 x}$.
\end{proof}

\begin{example}[Un problema de Cauchy, de principio a fin]\label{ex:b1:diffeq:cauchy2}
Resuélvase $y'' - 3y' + 2y = 0$ con $y(0) = 0$, $y'(0) = 1$. El
\omterm{def:b1:diffeq:linear2}{polinomio característico} $r^2 - 3r + 2 = (r - 1)(r - 2)$ tiene las
raíces reales $1$ y $2$: solución general $y = \lambda\eu^{x} +
\mu\eu^{2x}$. Las dos condiciones dan el sistema lineal
\[
\lambda + \mu = 0, \qquad \lambda + 2\mu = 1 ,
\]
luego $\mu = 1$, $\lambda = -1$:
\[
y(x) = \eu^{2x} - \eu^{x} .
\]
Comprobación: $y(0) = 0$; $y' = 2\eu^{2x} - \eu^x$ cumple $y'(0) = 1$;
y $y'' - 3y' + 2y = (4 - 6 + 2)\eu^{2x} + (-1 + 3 - 2)\eu^x = 0$.
Obsérvese la forma de la respuesta: cerca de $-\infty$ domina el modo
lento $-\eu^x$; cerca de $+\infty$, el modo rápido $\eu^{2x}$. Leer las
soluciones como superposiciones de modos con distintas tasas de
crecimiento o de decaimiento es el hábito rentable --- así se organiza
el reparto entre transitorio y régimen permanente del problema del fin
de semana.
\end{example}

\begin{omfigure}
\begin{tikzpicture}[xscale=1.35, yscale=1.55]
  \draw[->, thin] (0,0) -- (6.4,0) node[below] {$x$};
  \draw[->, thin] (0,-0.85) -- (0,1.25) node[left] {$y$};
  \draw[omDef, thick, domain=0:6.2, samples=400, smooth]
    plot (\x, {exp(-0.35*\x)*cos(4*\x r)});
  \draw[omThm, thick, domain=0:6.2, samples=120, smooth]
    plot (\x, {(1+1.8*\x)*exp(-1.1*\x)});
  \draw[omProp, thick, domain=0:6.2, samples=120, smooth]
    plot (\x, {1.15*exp(-0.45*\x) - 0.15*exp(-2.5*\x)});
  \node[omDef, font=\small] at (2.1,-0.62)
    {$\Delta < 0$: oscila};
  \node[omThm, font=\small] at (0.95,1.1) {$\Delta = 0$};
  \node[omProp, font=\small] at (4.9,0.38)
    {$\Delta > 0$: sin oscilación};
\end{tikzpicture}

{\small Los tres regímenes de $y'' + ay' + by = 0$ con soluciones que
decaen: oscilación amortiguada (raíces complejas), retorno crítico
(raíz doble) y decaimiento sobreamortiguado (dos raíces reales). Qué
régimen se da se lee solo del signo de $\Delta = a^2 - 4b$ --- antes de
resolver nada.}
\end{omfigure}

\begin{theorem}[Estructura y problema de Cauchy]\label{thm:b1:diffeq:structure2}
\begin{enumerate}
  \item Si $y_p$ es una solución particular de $(E)$, las soluciones
        de $(E)$ son exactamente $y_p + y_h$, con $y_h$ recorriendo
        las soluciones de $(H)$.
  \item (Superposición) Si $y_1$ resuelve $y'' + ay' + by = f_1$ e
        $y_2$ resuelve $y'' + ay' + by = f_2$, entonces $y_1 + y_2$
        resuelve la ecuación con miembro derecho $f_1 + f_2$.
  \item Para todos $x_0 \in I$ y $(y_0, y_0')$, el \omterm{thm:b1:diffeq:voc}{problema de Cauchy}
        «$(E)$, $y(x_0) = y_0$, $y'(x_0) = y_0'$» tiene exactamente una
        solución en $I$. \emph{(La existencia, dada una solución
        particular; la unicidad, en general.)}
\end{enumerate}
\end{theorem}

\begin{proof}
(1) $y$ resuelve $(E)$ si y solo si $y - y_p$ resuelve $(H)$, por
linealidad de $y \mapsto y'' + ay' + by$. (2) es la misma linealidad.

(3) Por (1) basta demostrar que las constantes $(\lambda, \mu)$ siempre
se pueden ajustar, de manera única, a cualesquiera datos
$(y_0, y_0')$. Trasladando la variable, supóngase $x_0 = 0$. En el caso
(1) del~\cref{thm:b1:diffeq:homogeneous2}, $y = \lambda\eu^{r_1x} +
\mu\eu^{r_2x}$ da
\[
y(0) = \lambda + \mu, \qquad y'(0) = r_1\lambda + r_2\mu :
\]
un sistema lineal en $(\lambda, \mu)$ cuyo determinante es
$r_2 - r_1 \neq 0$; al resolverlo explícitamente,
$\mu = \frac{y_0' - r_1y_0}{r_2 - r_1}$ y $\lambda = y_0 - \mu$:
exactamente una solución. En el caso (2), $y(0) = \lambda$ y
$y'(0) = r_0\lambda + \mu$: el sistema es triangular de determinante
$1$, resuelto por $\lambda = y_0$, $\mu = y_0' - r_0y_0$. En el caso
(3), $y(0) = \lambda$ e $y'(0) = \alpha\lambda + \omega\mu$:
determinante $\omega \neq 0$, resuelto por $\lambda = y_0$,
$\mu = \frac{y_0' - \alpha y_0}\omega$. En cada caso, la \omterm{def:b1:logic:map}{aplicación}
$(\lambda, \mu) \mapsto (y(x_0), y'(x_0))$ es una \omterm{def:b1:logic:inj}{biyección} lineal ---
el lenguaje del~\cref{ch:b1:linmaps} comprimirá esta discusión por
casos en una sola frase.
\end{proof}

\begin{method}[Solución particular para $f(x) = P(x)\,\eu^{\gamma x}$]\label{met:b1:diffeq:particular}
Cuando el miembro derecho es $P(x)\,\eu^{\gamma x}$ con $P$ polinomio y
$\gamma \in \R$ (lo que cubre polinomios, exponenciales y, con $\gamma$
complejo o por superposición, $\cos$ y $\sin$), búsquese una solución
particular de la forma
\[
y_p(x) = x^{m}\, Q(x)\, \eu^{\gamma x},
\qquad
m = \text{multiplicidad de } \gamma \text{ como raíz de } \chi
\ (m = 0, 1 \text{ o } 2),
\]
con $Q$ un polinomio del mismo grado que $P$, cuyos coeficientes se
hallan sustituyendo e identificando. Para $f = K\cos\omega x$ (o
$\sin$), resuélvase con miembro derecho $K\eu^{\iu\omega x}$ y tómese
la parte real (resp.\ imaginaria).
\end{method}

\begin{example}[La superposición en acción]\label{ex:b1:diffeq:superposition}
Resuélvase $y'' - y = \eu^{x} + 4$ en $\R$. Homogénea:
$\chi(r) = r^2 - 1$, raíces $\pm1$, luego
$y_h = \lambda\eu^x + \mu\eu^{-x}$. Sepárese el miembro derecho y
trátese cada trozo con el método. \emph{Trozo $\eu^x$:} aquí
$\gamma = 1$ es raíz simple de $\chi$, así que se prueba
$y_1 = c\,x\,\eu^x$: entonces
$y_1'' - y_1 = c(x + 2)\eu^x - cx\eu^x = 2c\,\eu^x$, de donde
$c = \frac12$. \emph{Trozo $4$:} $\gamma = 0$ no es raíz; sirve la
constante $y_2 = -4$. Por superposición
(\cref{thm:b1:diffeq:structure2}~(2)):
\[
y = \frac{x\,\eu^x}2 - 4 + \lambda\,\eu^x + \mu\,\eu^{-x},
\qquad (\lambda, \mu) \in \R^2 .
\]
Obsérvese cómo los dos trozos han exigido formas \emph{distintas}
($m = 1$ frente a $m = 0$): el criterio de multiplicidad se aplica a
cada exponente por separado, que es justamente la razón de separar el
miembro derecho antes de conjeturar.
\end{example}

\begin{example}[La regla de la multiplicidad en acción]\label{ex:b1:diffeq:multiplicity}
Resuélvase $y'' + y' = x$ en $\R$. El miembro derecho es
$P(x)\eu^{0 \cdot x}$ con $P(x) = x$, y $\gamma = 0$ es raíz
\emph{simple} de $\chi(r) = r^2 + r = r(r + 1)$: luego $m = 1$, y la
conjetura correcta es $y_p = x\,(\alpha x + \beta) = \alpha x^2 +
\beta x$, un grado más que $P$. Sustituyendo:
\[
y_p'' + y_p' = 2\alpha + (2\alpha x + \beta)
= 2\alpha x + (2\alpha + \beta) ,
\]
e identificando con $x$ se obtiene $\alpha = \frac12$, $\beta = -1$:
$y_p = \frac{x^2}2 - x$. Solución general:
$y = \frac{x^2}2 - x + \lambda + \mu\,\eu^{-x}$. Si se hubiese
conjeturado $y_p = \alpha x + \beta$ (ignorando la multiplicidad), la
sustitución daría $y_p'' + y_p' = \alpha$, una constante --- ninguna
elección de $\alpha, \beta$ puede igualar a $x$, y el fracaso es
estructural: las constantes ya resuelven la ecuación homogénea, así que
son invisibles para el miembro izquierdo. El factor $x^m$ existe
precisamente para salir del espacio de soluciones homogéneas.
\end{example}

\begin{remark}[La póliza de seguros de treinta segundos]\label{rem:b1:diffeq:check}
Toda ecuación resuelta en este capítulo termina con una comprobación
por sustitución, y no es un adorno. Un cálculo con ecuaciones
diferenciales encadena muchos pasos pequeños (una primitiva, una regla
del producto, dos constantes), y un solo error de signo se propaga
invisible; sustituir la fórmula final en la ecuación los caza casi
todos al precio de una derivación. Cultívese el reflejo en tres capas:
compruébese la \emph{solución particular} sola (la parte homogénea se
cancela de todos modos), compruébense las \emph{condiciones iniciales}
en la solución completa y, cuando haya un parámetro, compruébese un
\emph{valor degenerado} (¿reproduce la fórmula para $\Omega$ general la
respuesta conocida en $\Omega = 0$?). El hábito cuesta medio minuto y
convierte «probablemente correcto» en «verificado».
\end{remark}

\begin{remark}[Errores frecuentes]\label{rem:b1:diffeq:pitfalls}
\begin{enumerate}
  \item \emph{Normalícese primero.} Las fórmulas suponen que la
        ecuación se lee $y' + a(x)y = b(x)$ --- coeficiente $1$ en
        $y'$. Para $xy' - 2y = x^3$, divídase por $x$ (en un intervalo
        que evite $0$) antes de identificar $a$ y $b$, como en
        el~\cref{exo:b1:diffeq:2}.
  \item \emph{Una constante por dimensión, fijada al final.} La
        solución general de primer orden lleva una constante; la de
        segundo orden, dos; las condiciones iniciales se imponen sobre
        la solución \emph{completa} $y_p + y_h$, nunca sobre $y_h$
        sola --- imponerlas antes de sumar $y_p$ es el error
        estructural más frecuente.
  \item \emph{Ojo con la multiplicidad.} Una conjetura de solución
        particular que resuelva la ecuación homogénea es invisible
        para el miembro izquierdo; el factor $x^m$
        del~\cref{met:b1:diffeq:particular} no es opcional
        (\cref{ex:b1:diffeq:multiplicity}).
  \item \emph{Los intervalos forman parte de la respuesta.} Las
        soluciones viven en intervalos donde los coeficientes son
        continuos; pegar a través de una singularidad puede crear
        constantes espurias (\cref{exo:b1:diffeq:12}) o destruir la
        unicidad. «Resuélvase en $\R^*$» significa dos problemas
        independientes.
\end{enumerate}
\end{remark}

\begin{example}[Excitación fuera de resonancia]\label{ex:b1:diffeq:offresonance}
Resuélvase $y'' + 4y = \sin x$ en $\R$. Frecuencia propia $2$,
frecuencia de excitación $1$: como $\iu$ \emph{no} es raíz de
$\chi(r) = r^2 + 4$, la multiplicidad es $m = 0$ y basta con una
sinusoide simple. Probando $y_p = \alpha\sin x$ (no hace falta coseno:
la ecuación no tiene término en $y'$ y $\sin$ regenera $\sin$):
\[
y_p'' + 4y_p = -\alpha\sin x + 4\alpha\sin x = 3\alpha\sin x ,
\]
luego $\alpha = \frac13$ y la solución general es
\[
y = \frac{\sin x}3 + \lambda\cos 2x + \mu\sin 2x .
\]
Toda solución permanece acotada: una superposición de dos oscilaciones,
de frecuencias $1$ (forzada) y $2$ (propia). Compárese con el ejemplo
siguiente, donde excitar \emph{a} la frecuencia propia cambia la forma
misma de la respuesta.
\end{example}

\begin{example}[Una oscilación forzada]\label{ex:b1:diffeq:oscillation}
Resuélvase $y'' + y = \cos x$, $y(0) = 0$, $y'(0) = 0$.

\emph{Homogénea:} $\chi(r) = r^2 + 1$, raíces $\pm\iu$:
$y_h = \lambda\cos x + \mu \sin x$.

\emph{Particular:} miembro derecho $\Re(\eu^{\iu x})$ con
$\gamma = \iu$ raíz simple de $\chi$: pruébese $z_p = c\, x\,
\eu^{\iu x}$ ($c \in \C$). Entonces $z_p'' + z_p = c\,(2\iu)\eu^{\iu x}$,
que vale $\eu^{\iu x}$ para $c = \frac{1}{2\iu} = -\frac\iu2$. Luego
$z_p = -\frac{\iu}{2} x (\cos x + \iu \sin x)$ e
$y_p = \Re(z_p) = \frac{x \sin x}{2}$.

\emph{Solución general:} $y = \frac{x\sin x}{2} + \lambda\cos x +
\mu\sin x$. Condiciones: $y(0) = \lambda = 0$; $y' = \frac{\sin x +
x\cos x}{2} + \mu\cos x$, luego $y'(0) = \mu = 0$. Respuesta:
$y = \frac{x\sin x}{2}$ --- una oscilación cuya amplitud crece
linealmente: el fenómeno de \emph{resonancia}\index{resonancia},
causado por excitar el sistema a su frecuencia propia.
\end{example}

\begin{omfigure}
\begin{tikzpicture}
\begin{axis}[omaxis, width=11cm, height=5.2cm,
    xmin=0, xmax=25, ymin=-13, ymax=13,
    xtick={0,5,10,15,20,25}, ytick=\empty]
  \addplot[omThm, thick, domain=0:25, samples=400] {x*sin(deg(x))/2};
  \addplot[dashed, gray, domain=0:25] {x/2};
  \addplot[dashed, gray, domain=0:25] {-x/2};
\end{axis}
\end{tikzpicture}

{\small \omterm{ex:b1:diffeq:oscillation}{Resonancia}: la solución $y = \frac{x \sin x}{2}$ de
$y'' + y = \cos x$ oscila entre las rectas $y = \pm\frac x2$ (a
trazos), con amplitud siempre creciente.}
\end{omfigure}

\begin{remark}[Interludio: la linealidad es una geometría]\label{rem:b1:diffeq:interlude}
Obsérvese la forma de todos los \omterm{def:b1:logic:sets}{conjuntos} de soluciones de este
capítulo: una solución especial más un espacio de soluciones homogéneas
con una constante libre (primer orden) o dos (segundo orden). Los
capítulos de álgebra lineal
(\cref{ch:b1:vspaces,ch:b1:findim,ch:b1:linmaps}) aportarán el
vocabulario exacto: la \omterm{def:b1:logic:map}{aplicación} $L(y) = y'' + ay' + by$ es
\emph{lineal}, sus soluciones homogéneas forman el \emph{núcleo} de
$L$, un espacio vectorial cuya \emph{dimensión} es el orden de la
ecuación --- ese es el contenido honesto de «una constante por orden»
--- y el \omterm{def:b1:logic:sets}{conjunto} de soluciones de $L(y) = f$ es un \emph{subespacio
afín}, un trasladado del núcleo. Incluso la \omterm{def:b1:logic:map}{aplicación} de Cauchy
$(\lambda, \mu) \mapsto (y(x_0), y'(x_0))$
del~\cref{thm:b1:diffeq:structure2} es una \omterm{def:b1:logic:inj}{biyección} lineal entre dos
planos, es decir, un sistema $2 \times 2$ invertible
(\cref{ch:b1:matrices}). Nada de este capítulo habrá que rehacerlo ---
solo habrá que rebautizarlo, y ese rebautizo es el mejor
calentamiento posible para el álgebra lineal: allí, toda definición
abstracta ya se ha ganado la vida aquí.
\end{remark}

\begin{remark}[Dónde se usa este capítulo]\label{rem:b1:diffeq:whereused}
El teorema de estructura --- las soluciones de $(E)$ son «una solución
particular más las soluciones de $(H)$» --- es la primera aparición de
un patrón que los \cref{ch:b1:vspaces,ch:b1:linmaps} bautizarán: el
\omterm{def:b1:logic:sets}{conjunto} de soluciones de $(H)$ es el \emph{núcleo} de la
\omterm{def:b1:logic:map}{aplicación} lineal $y \mapsto y'' + ay' + by$, y el \omterm{def:b1:logic:sets}{conjunto} de
soluciones de $(E)$ es un trasladado afín suyo. El
\omterm{def:b1:diffeq:linear2}{polinomio característico} reaparece como \omterm{def:b1:diffeq:linear2}{polinomio característico}
de una matriz en el~\cref{ch:b1:matrices}: una ecuación de segundo
orden es un sistema de primer orden $2 \times 2$ disfrazado, punto de
vista que el volumen del segundo año sistematiza. Las integrales que
exige la \omterm{thm:b1:diffeq:voc}{variación de las constantes} las suministra
el~\cref{ch:b1:integration}, y el problema del fin de semana --- el
oscilador amortiguado y forzado --- es el caso modelo de toda cuestión
de oscilaciones en las ciencias, de los circuitos a los puentes
colgantes.
\end{remark}

\section{Ejercicios}

\begin{exercise}[$\star$]\label{exo:b1:diffeq:1}
Resuélvase en $\R$: $\;y' + 2y = \eu^{3x}$; después el
\omterm{thm:b1:diffeq:voc}{problema de Cauchy} $y(0) = 1$.
\end{exercise}

\begin{exercise}[$\star$]\label{exo:b1:diffeq:2}
Resuélvase en $\intoo{0}{+\infty}$: $\;x y' - 2y = x^3$
\emph{(póngase primero la ecuación en forma normalizada)}.
\end{exercise}

\begin{exercise}[$\star$]\label{exo:b1:diffeq:3}
Resuélvanse en $\R$, dando la solución general real:
$\;y'' - 3y' + 2y = 0$; $\;y'' + 4y' + 4y = 0$; $\;y'' - 2y' + 5y =
0$.
\end{exercise}

\begin{exercise}[$\star$]\label{exo:b1:diffeq:4}
Resuélvase $y'' - y = x^2$ en $\R$, y después el
\omterm{thm:b1:diffeq:voc}{problema de Cauchy} $y(0) = 0$, $y'(0) = 1$.
\end{exercise}

\begin{exercise}[$\star\star$]\label{exo:b1:diffeq:5}
Resuélvase en $\intoo{-\frac\pi2}{\frac\pi2}$:
$\;y' + y\tan x = \sin 2x$.
\end{exercise}

\begin{exercise}[$\star\star$]\label{exo:b1:diffeq:6}
Resuélvase $y'' - 4y' + 3y = (2x + 1)\,\eu^{x}$ en $\R$. \emph{(Ojo con
la multiplicidad: ¿es $1$ raíz del \omterm{def:b1:diffeq:linear2}{polinomio característico}?)}
\end{exercise}

\begin{exercise}[$\star\star$]\label{exo:b1:diffeq:7}
Resuélvase $y'' + 4y = \sin 2x + x$ en $\R$ \emph{(superposición;
trátese cada miembro derecho por separado)}.
\end{exercise}

\begin{exercise}[$\star\star$]\label{exo:b1:diffeq:8}
Una taza de café a temperatura $T_0 = 80\,^\circ$C está en una
habitación a $20\,^\circ$C. La ley de enfriamiento de Newton dice que
$T' = -k\,(T - 20)$ con $k > 0$. Resuélvase para $T(t)$ y, sabiendo que
el café está a $50\,^\circ$C al cabo de $10$ minutos, hállese cuándo
llega a $25\,^\circ$C.
\end{exercise}

\begin{exercise}[$\star\star\star$]\label{exo:b1:diffeq:9}
(Oscilador amortiguado) Para $\varepsilon \geq 0$, considérese
$y'' + 2\varepsilon y' + y = 0$.
\begin{enumerate}
  \item Resuélvase para $\varepsilon \in \intco{0}{1}$,
        $\varepsilon = 1$ y $\varepsilon > 1$.
  \item Pruébese que, para todo $\varepsilon > 0$, todas las
        soluciones tienden a $0$ en $+\infty$, y que para
        $\varepsilon = 0$ las soluciones no nulas no lo hacen.
  \item Para $\varepsilon \in \intoo{0}{1}$, pruébese que los ceros de
        una solución no nula están regularmente espaciados, con paso
        $\frac{\pi}{\sqrt{1 - \varepsilon^2}}$.
\end{enumerate}
\end{exercise}

\begin{exercise}[$\star\star\star$]\label{exo:b1:diffeq:10}
Hállense todas las funciones $f \colon \R \to \R$, dos veces
derivables, tales que
\[
\forall x, y \in \R, \qquad f(x + y) + f(x - y) = 2 f(x) f(y),
\]
con $f(0) \ne 0$ y $f$ no constante.
\emph{Indicación: fíjese $y$ y derívese dos veces respecto de $x$ en
$0$; pruébese que $f(0) = 1$ y que $f'' = c f$ para alguna constante
$c$; resuélvase después según el signo de $c$ y compruébese qué
soluciones satisfacen la ecuación funcional.}
\end{exercise}

\begin{exercise}[$\star\star$]\label{exo:b1:diffeq:11}
(Ecuación de Euler) Resuélvase $x^2 y'' - x y' + y = 0$ en
$\intoo{0}{+\infty}$. \emph{Indicación: póngase $z(t) = y(\eu^t)$, es
decir, sustitúyase $x = \eu^t$, y véase que $z$ cumple una ecuación
lineal de coeficientes constantes.}
\end{exercise}

\begin{exercise}[$\star\star\star$]\label{exo:b1:diffeq:12}
Considérese la ecuación $x\,y' = 2y$ en toda la recta real, en la
incógnita, una función derivable $y \colon \R \to \R$.
\begin{enumerate}
  \item Resuélvase en $\intoo{0}{+\infty}$ y en $\intoo{-\infty}{0}$.
  \item Pruébese que, para constantes $a, b \in \R$
        \emph{cualesquiera}, la función que vale $ax^2$ para $x \geq 0$
        y $bx^2$ para $x < 0$ es derivable en $\R$ y resuelve la
        ecuación en todas partes.
  \item Conclúyase que el \omterm{def:b1:logic:sets}{conjunto} de soluciones en $\R$ es una
        familia de dos parámetros, y explíquese por qué esto no
        contradice la unicidad del~\cref{thm:b1:diffeq:voc}.
\end{enumerate}
\end{exercise}

\section{Problema: el oscilador amortiguado y forzado}

\begin{problem}\label{pb:b1:diffeq:1}
Una sola ecuación gobierna una masa unida a un muelle en un medio
viscoso, la carga de un circuito RLC y un edificio que oscila con el
viento:
\[
(E_\Omega)\colon\quad
x'' + 2\lambda x' + \omega_0^2\,x = A\cos(\Omega t),
\]
con $\lambda \geq 0$ el amortiguamiento, $\omega_0 > 0$ la frecuencia
propia y $A > 0$, $\Omega > 0$ la amplitud y la frecuencia de la
excitación. Este problema extrae su comportamiento completo: el
decaimiento de los transitorios, el único régimen permanente periódico,
la curva de \omterm{ex:b1:diffeq:oscillation}{resonancia} y su agudeza (el \emph{\omterm{pb:b1:diffeq:1}{factor de calidad}}),
las \omterm{pb:b1:diffeq:1}{pulsaciones} del caso sin amortiguamiento y el balance energético
que sostiene la oscilación.\index{resonancia}\index{oscilador}
Salvo indicación en contra, $0 < \lambda < \omega_0$ (régimen
subamortiguado) y se escribe $\omega_d = \sqrt{\omega_0^2 - \lambda^2}$.

\medskip
\noindent\textbf{Parte I --- El oscilador libre.} Aquí $A = 0$.
\begin{enumerate}
  \item Resuélvase la ecuación homogénea $(H)$ para
        $0 < \lambda < \omega_0$ y para $\lambda = 0$. (Los regímenes
        $\lambda \geq \omega_0$ se trataron en
        el~\cref{exo:b1:diffeq:9}; cítense.)
  \item Pruébese que, para todo $\lambda > 0$, todas las soluciones de
        $(H)$ tienden a $0$ en $+\infty$ --- en los tres regímenes.
  \item Defínase la energía $\mathcal E(t) = \frac12 x'(t)^2 +
        \frac12\omega_0^2\,x(t)^2$ a lo largo de una solución de $(H)$.
        Pruébese que $\mathcal E'(t) = -2\lambda\,x'(t)^2 \leq 0$ y
        dedúzcase (sin resolver nada) que el \omterm{thm:b1:diffeq:voc}{problema de Cauchy}
        «$(H)$, $x(t_0) = x'(t_0) = 0$» solo tiene la solución nula,
        para todo $\lambda \geq 0$.
  \item Para $0 < \lambda < \omega_0$, escríbase la solución no nula
        como $x(t) = R\,\eu^{-\lambda t}\cos(\omega_d t - \varphi)$ y
        sea $T_d = \frac{2\pi}{\omega_d}$ el seudoperíodo. Pruébese que
        $x(t + T_d) = \eu^{-\lambda T_d}\,x(t)$: cada oscilación es la
        anterior encogida por el factor constante $\eu^{-\delta}$, con
        $\delta = \frac{2\pi\lambda}{\omega_d}$ (el \emph{decremento
        logarítmico}). Calcúlese $\delta$ para $\omega_0 = 1$,
        $\lambda = 0.1$.
  \item Defínase el \emph{\omterm{pb:b1:diffeq:1}{factor de calidad}} $Q =
        \dfrac{\omega_0}{2\lambda}$\index{factor de calidad}. Pruébese
        que, tras el tiempo $\frac1\lambda$ (una división de la
        amplitud por $\eu$), el oscilador ha completado
        $\frac{\omega_d}{2\pi\lambda}$ seudoperíodos, que para
        amortiguamiento débil ($\lambda \ll \omega_0$) es
        aproximadamente $\frac Q\pi$: el \omterm{pb:b1:diffeq:1}{factor de calidad} cuenta,
        salvo el factor $\pi$, las oscilaciones que sobreviven antes de
        que la amplitud se divida por $\eu$.
\end{enumerate}

\medskip
\noindent\textbf{Parte II --- El régimen permanente.} Ahora $A > 0$ y
$\lambda > 0$.
\begin{enumerate}[resume]
  \item Búsquese una solución particular como parte real de
        $z\,\eu^{\iu\Omega t}$ con $z \in \C$
        (\cref{met:b1:diffeq:particular}). Pruébese que funciona con
        \[
        z = \frac{A}{\omega_0^2 - \Omega^2 + 2\iu\lambda\Omega} .
        \]
  \item Dedúzcase el régimen permanente en forma amplitud--fase:
        $x_p(t) = R(\Omega)\cos\bigl(\Omega t -
        \varphi(\Omega)\bigr)$ con
        \[
        R(\Omega) = \frac{A}{\sqrt{(\omega_0^2 - \Omega^2)^2 +
        4\lambda^2\Omega^2}},
        \qquad
        \tan\varphi = \frac{2\lambda\Omega}{\omega_0^2 -
        \Omega^2},
        \quad \varphi \in \intoo0\pi .
        \]
  \item Interprétense los dos regímenes extremos: calcúlense los
        límites de $R$ y de $\varphi$ cuando $\Omega \to 0^+$
        (respuesta cuasiestática $A/\omega_0^2$, fase $0$) y cuando
        $\Omega \to +\infty$ ($R \sim A/\Omega^2 \to 0$, fase
        $\to \pi$: la masa se mueve en sentido opuesto a una excitación
        demasiado rápida).
  \item Pruébese que \emph{toda} solución de $(E_\Omega)$ es $x_p$ más
        una solución de $(H)$ y que, por tanto, converge al régimen
        permanente $x_p$ cuando $t \to +\infty$, sean cuales sean las
        condiciones iniciales: una vez muerto el transitorio, el
        oscilador no recuerda cómo empezó.
  \item Pruébese que $x_p$ es la \emph{única} solución periódica de
        $(E_\Omega)$.
  \item Resuélvase un \omterm{thm:b1:diffeq:voc}{problema de Cauchy} hasta el final: para
        $x'' + 2x' + 2x = \cos t$ con $x(0) = x'(0) = 0$, pruébese que
        la solución es
        \[
        x(t) = \frac{\cos t + 2\sin t}5
        - \eu^{-t}\,\frac{\cos t + 3\sin t}5 ,
        \]
        e identifíquense las partes transitoria y permanente.
\end{enumerate}

\medskip
\noindent\textbf{Parte III --- La curva de \omterm{ex:b1:diffeq:oscillation}{resonancia}.} Estudio de
$\Omega \mapsto R(\Omega)$ en $\intoo0{+\infty}$.
\begin{enumerate}[resume]
  \item Poniendo $u = \Omega^2$ y $g(u) = (\omega_0^2 - u)^2 +
        4\lambda^2 u$, pruébese que, si $2\lambda^2 < \omega_0^2$,
        entonces $R$ alcanza un máximo estricto en la
        \emph{frecuencia de \omterm{ex:b1:diffeq:oscillation}{resonancia}} $\Omega_r =
        \sqrt{\omega_0^2 - 2\lambda^2}$, con
        \[
        R_{\max} = R(\Omega_r)
        = \frac{A}{2\lambda\sqrt{\omega_0^2 - \lambda^2}} .
        \]
  \item Pruébese que $\dfrac{R_{\max}}{R(0)} = Q\,\bigl(1 -
        \tfrac{\lambda^2}{\omega_0^2}\bigr)^{-1/2} \geq Q$: en
        \omterm{ex:b1:diffeq:oscillation}{resonancia}, la excitación se amplifica (esencialmente) por
        el \omterm{pb:b1:diffeq:1}{factor de calidad}.
  \item Demuéstrese que la amplitud de la \emph{velocidad}
        $V(\Omega) = \Omega\,R(\Omega)$ es máxima exactamente en
        $\Omega = \omega_0$ (no en $\Omega_r$), y que la fase ahí vale
        $\varphi(\omega_0) = \frac\pi2$: en $\Omega = \omega_0$, la
        velocidad está exactamente en fase con la fuerza.
  \item (Ancho de banda) Resuélvase exactamente $g(u) = 2\,g(u_r)$,
        donde $u_r = \omega_0^2 - 2\lambda^2$, y dedúzcase que las dos
        frecuencias $\Omega_\pm$ en las que $R = R_{\max}/\sqrt2$
        cumplen $\Omega_+^2 - \Omega_-^2 = 4\lambda\sqrt{\omega_0^2 -
        \lambda^2}$; conclúyase que, para amortiguamiento débil, el
        ancho de banda es $\Omega_+ - \Omega_- \approx 2\lambda$, es
        decir, $Q \approx \frac{\omega_0}{\Omega_+ - \Omega_-}$: los
        picos de \omterm{ex:b1:diffeq:oscillation}{resonancia} agudos son sistemas de $Q$ alto.
  \item Retrato numérico para $\omega_0 = 1$, $\lambda = 0.05$
        ($Q = 10$), $A = 1$: calcúlense $\Omega_r$, $R_{\max}$, la
        respuesta estática $R(0)$ y el ancho de banda aproximado.
  \item Pruébese que, si $2\lambda^2 \geq \omega_0^2$, entonces $R$ es
        estrictamente decreciente en $\intoo0{+\infty}$: los sistemas
        muy amortiguados no tienen pico de \omterm{ex:b1:diffeq:oscillation}{resonancia} alguno.
\end{enumerate}

\medskip
\noindent\textbf{Parte IV --- Sin amortiguamiento: \omterm{pb:b1:diffeq:1}{pulsaciones} y
\omterm{ex:b1:diffeq:oscillation}{resonancia}.} Aquí $\lambda = 0$.
\begin{enumerate}[resume]
  \item Para $\Omega \neq \omega_0$, hállese la solución general de
        $x'' + \omega_0^2 x = A\cos(\Omega t)$.
  \item Resuélvase el \omterm{thm:b1:diffeq:voc}{problema de Cauchy} $x(0) = x'(0) = 0$ y
        transfórmese la respuesta en la forma de producto
        \[
        x(t) = \frac{2A}{\omega_0^2 - \Omega^2}\,
        \sin\Bigl(\frac{(\omega_0 - \Omega)t}2\Bigr)
        \sin\Bigl(\frac{(\omega_0 + \Omega)t}2\Bigr) .
        \]
  \item Para $\Omega$ próximo a $\omega_0$, léase el producto como una
        oscilación rápida de frecuencia $\frac{\omega_0 + \Omega}2$
        modulada por una envolvente lenta de frecuencia
        $\frac{\abs{\omega_0 - \Omega}}2$: las \emph{pulsaciones}
        \index{pulsaciones}. Dense el período de la envolvente y la
        amplitud máxima, y obsérvese cómo ambos estallan cuando
        $\Omega \to \omega_0$.
  \item Fíjese $t$ y hágase $\Omega \to \omega_0$ en la fórmula de la
        pregunta 19: pruébese que el límite es
        \[
        x_\infty(t) = \frac{A\,t\,\sin(\omega_0 t)}{2\omega_0} ,
        \]
        y compruébese directamente que $x_\infty$ resuelve la ecuación
        resonante $x'' + \omega_0^2 x = A\cos(\omega_0 t)$ con
        $x(0) = x'(0) = 0$ (compárese con
        el~\cref{ex:b1:diffeq:oscillation}): la \omterm{ex:b1:diffeq:oscillation}{resonancia} es el límite
        de \omterm{pb:b1:diffeq:1}{pulsaciones} cada vez más lentas y más amplias.
  \item Contrástense los dos destinos de la \omterm{ex:b1:diffeq:oscillation}{resonancia}: crecimiento
        lineal $\frac{At}{2\omega_0}$ sin amortiguamiento frente a
        saturación en $R_{\max} \approx Q\,\frac{A}{\omega_0^2}$ con
        amortiguamiento débil. En una frase: ¿qué mecanismo físico
        convierte el primero en la segunda?
\end{enumerate}

\medskip
\noindent\textbf{Parte V --- Balance energético y síntesis.}
\begin{enumerate}[resume]
  \item En el régimen permanente de la parte II, calcúlese el
        promedio en un período $\frac{2\pi}\Omega$ de (a) la potencia
        aportada por la excitación, $P_{\mathrm{in}}(t) =
        A\cos(\Omega t) \cdot x_p'(t)$, y (b) la potencia disipada por
        el amortiguamiento, $P_{\mathrm{diss}}(t) =
        2\lambda\,x_p'(t)^2$. Pruébese que los dos promedios valen
        $\lambda\,R^2\Omega^2$: la excitación aporta exactamente lo que
        el amortiguamiento quema --- por eso el régimen permanente es
        permanente.
  \item ¿Dónde ha usado exactamente el problema: (i) el teorema de
        estructura \cref{thm:b1:diffeq:structure2}; (ii) el método de
        la exponencial compleja; (iii) un estudio de función de
        variable real al estilo del~\cref{ch:b1:functions}? Una frase
        para cada uno.
  \item Síntesis: descríbase el mapa completo de comportamientos de
        $(E_\Omega)$ --- libre frente a forzado, amortiguado frente a
        no amortiguado, el papel de $Q$ como único mando adimensional
        que ajusta altura del pico, ancho de banda y vida del
        transitorio --- y menciónese dónde continúa la historia:
        sistemas de primer orden $2 \times 2$ (\cref{ch:b1:matrices} y
        el volumen del segundo año) y la descomposición de una
        excitación periódica general en sinusoides (series de Fourier,
        en el volumen del tercer año), para la cual el caso sinusoidal
        de este problema es el ladrillo fundamental.
\end{enumerate}
\end{problem}
